\specialsection{Заключение.} \label{conclusion2}

В рамках работы создана библиотека Compio, решающая актуальную задачу эффективного хранения и модификации сжатых данных с произвольным доступом. В отличие от существующих аналогов, библиотека поддерживает не только чтение, но и полноценную запись, включая вставку и удаление фрагментов в середине логического файла, без необходимости полной перепаковки данных.

Ключевыми особенностями разработанного решения являются:
\begin{itemize}
    \item универсальность — возможность подключения произвольных алгоритмов сжатия;
    \item самодостаточность формата файла-контейнера, не требующего внешних индексов;
    \item эффективное управление внутренней фрагментацией блоков, подтверждённое экспериментально;
    \item двухуровневое кэширование, минимизирующее операции ввода-вывода и сжатия;
    \item совместимость с любыми языками программирования через стабильный C ABI.
\end{itemize}

Практическая значимость библиотеки заключается в её применимости в областях, где необходимо хранить большие объёмы данных в сжатом виде и при этом регулярно выполнять операции произвольной записи или редактирования, — биоинформатика, астрономия, системы логирования и аналитики.

Дальнейшие направления развития включают разработку обёрток для языков высокого уровня (в первую очередь Python), экспериментальную проверку работы библиотеки на файлах существенно большего объёма (десятки и сотни гигабайт) для оценки масштабируемости алгоритмов, реализацию многопоточной обработки, использующей независимость сжатия отдельных блоков для дополнительного ускорения операций ввода-вывода, а также исследование возможности поиска непосредственно в сжатых данных.
